%! TEX root = thesis.tex

\chapter{Related Work}
\label{sec:related_work}

\section{Temporal Workload Shifting}

Temporal workload shifting delays computational jobs to align with periods of
low grid carbon intensity.  Wiesner~et~al.~\cite{wiesner_lets_2021}
demonstrated that delaying cloud batch jobs by a few hours can reduce
operational emissions by 3--34\%, with negligible checkpoint overhead enabling
higher savings.  Their work uses a single threshold for pause decisions and
does not explore hysteresis.  Carbon Explorer~\cite{acun_carbon_2023} balances
operational and embodied carbon at the datacenter planning timescale, but
operates at the level of hardware procurement and deployment decisions rather
than per-run scheduling.

\section{Carbon-Aware Distributed Training}

Extending temporal shifting to geo-distributed training adds spatial
flexibility.  Wiesner~et~al.~\cite{wiesner_distributed_2026} trained a
561M-parameter transformer across multiple sites during renewable curtailment
windows, reducing emissions to 5--12\% of the baseline.  Their approach
exploits near-zero-carbon curtailment energy but requires multiple
geographically distributed data centers.  Our work complements this by
optimizing the hysteresis policy for single-site training, which is the more
common deployment scenario.

\section{Emissions Accounting Methodology}

Accurate carbon accounting is essential for evaluating the savings from
temporal shifting.  Patterson~et~al.~\cite{patterson_carbon_nodate} propose a
methodology based on GPU-hours, thermal design power, power usage effectiveness, and grid intensity.
Dodge~et~al.~\cite{dodge_measuring_2022} highlight the importance of using
marginal rather than average grid intensity, and demonstrate that
time-specific cloud data yields more accurate estimates.  We follow their
recommendation by using marginal carbon intensity at 5-minute resolution.


